\chapter{Overview}
\label{ch:overview}

The \productname\ is an ultrasonic wedge wire bonder. It joins fine wire to
metallised pads by pressing the wire down with a controlled force and vibrating
it ultrasonically until the two metals weld. No heat and no solder are
involved: the bond is made by friction and pressure alone.

A complete cycle makes two bonds and the loop of wire between them, then tears
the wire and forms a fresh tail ready for the next cycle.

\figureplaceholder{Photograph of the complete machine with the bond head, Y
stage, control panel, display and mouse labelled.}

\section{The axes}

Three axes position the tool and the work.

\begin{center}\small
\begin{tabular}{@{}L{16mm}L{34mm}L{72mm}@{}}
\toprule
\textbf{Axis} & \textbf{What it moves} & \textbf{Notes} \\
\midrule
Z & The bond head, vertically &
  9\,mm of travel. Position is measured absolutely, so the machine knows the
  head's height the instant it powers up --- Z never needs homing. \\
\addlinespace
Y & The workholder stage &
  18\,mm of travel. Homed against a limit switch at power-up; until that
  succeeds nothing may move. \\
\addlinespace
T & The wire feed &
  Feeds the tail and tears the wire. It has no home switch, so every T move is
  relative to wherever the axis currently sits. \\
\bottomrule
\end{tabular}
\end{center}

\begin{wbnote}
Because the T axis is relative, \menuitem{Tear} is applied on top of wherever
\menuitem{Tail} left the axis. The two are not independent coordinates ---
changing \menuitem{Tail} moves the tear endpoint too.
\end{wbnote}

\section{The bonding head}

\begin{description}[leftmargin=!,labelwidth=32mm,style=nextline,font=\normalfont\bfseries]

\item[Force coil] Produces the bond force. You set force in grams; the machine
  converts that to a coil current using a calibration that you can correct from
  the panel (Chapter~\ref{ch:calibration}). Force is applied smoothly rather
  than stepped, so releasing it after a weld does not flick the bond lever.

\item[Ultrasonic transducer] Vibrates the tool. Before each weld the machine
  sweeps a band of frequencies to find where the transducer resonates, then
  drives it there and counts the energy going in until the requested dose is
  delivered.

\item[Wire clamp] Holds or releases the wire. It opens to let wire pay out
  while the loop forms, and closes to hold the wire for the tear.

\item[Contact sensor] Detects electrical contact between tool and workpiece.
  The machine uses it to know when it has touched down, and when the wire has
  released on the way back up.

\end{description}

\figureplaceholder[0.45\textwidth]{Close-up of the bond head showing the tool,
transducer, wire clamp, wire path and force lever.}

\section{How a cycle works}

Every bonding mode follows the same outline. What differs is how the descents
are commanded and how the wire is torn.

\begin{enumerate}[leftmargin=*]
\item You press the mouse button. The head descends and searches for the first
  pad.
\item It touches down, settles, ramps to the first bond force, and welds.
\item The clamp opens and the head rises, kinking the wire and feeding a tail.
\item The head rises to the loop apex and waits for you.
\item You press again. The head moves across, descends and welds the second
  bond.
\item The clamp closes and the wire is torn.
\item A fresh tail is formed and the machine returns to its parked height.
\end{enumerate}

Chapter~\ref{ch:modes} covers each mode's version of this in detail.

\section{Before every cycle}

Whatever the mode, the machine always puts itself into the same known state
before it moves: the force coil is switched off, the clamp is closed, the head
is driven to \menuitem{Reset Height}, and any button presses left over from
before are discarded.

\begin{wbnote}
That last point matters in practice: a button you pressed while the machine was
busy cannot trigger a cycle later. Every cycle starts from a deliberate press.
\end{wbnote}
